\chapter[Transformata Fouriera i dystrybucje temperowane]{Transformata Fouriera, przestrzeń Schwartza i dystrybucje temperowane}
\label{chap:pde-fourier-Rd}

Rozdział~\ref{chap:pde-distributions} pozwolił różniczkować funkcje
niegładkie i rozwiązywać równania o stałych współczynnikach za pomocą
rozwiązań fundamentalnych. Transformata Fouriera realizuje komplementarny
pomysł: zamienia pochodne na mnożenie przez wielomiany. Na całej
$\R^d$ nie możemy jednak bez zastrzeżeń używać klasycznej transformaty
całkowej dla wszystkich dystrybucji. Potrzebujemy więc przestrzeni
Schwartza $\mathcal S(\R^d)$ i jej przestrzeni sprzężonej $\mathcal S'(\R^d)$.

W całym rozdziale $d\ge1$, całki na $\R^d$ są liczone względem
miary Lebesgue'a $\lambda_d$, a wielowskaźniki należą do $\N_0^d$.
Dla funkcji całkowalnych przyjmujemy \emph{nieunitarną}
konwencję
\begin{equation}\label{eq:pde8-convention}
 \widehat f(\xi)=\mathcal F f(\xi)
 =\int_{\R^d}e^{-ix\cdot\xi}f(x)\,\mathrm d\lambda_d(x),\qquad
 \mathcal F^{-1}g(x)=(2\pi)^{-d}
 \int_{\R^d}e^{ix\cdot\xi}g(\xi)\,\mathrm d\lambda_d(\xi).
\end{equation}
Na razie symbol $\mathcal F^{-1}$ oznacza podaną całkę
dla $g\in L^1(\R^d)$; fakt, że jest to odwzorowanie
odwrotne do $\mathcal F$ na przestrzeni Schwartza,
wykażemy później.
Współczynniki $(2\pi)^{-d}$ zależą od konwencji i nie będą
pomijane. W szczególności nie będziemy bez przeliczenia używać
wzorów wyprowadzonych dla transformaty unitarnej.
Nie należy też utożsamiać tej transformaty całkowej na $\R^d$
z dyskretnymi współczynnikami szeregów Fouriera z
rozdziału~\ref{chap:pde-fourier}: te ostatnie są uśredniane
przez czynnik $(2\pi)^{-d}$.

\section{Przestrzeń Schwartza: po co wymagać szybkiego zaniku?}
\label{sec:pde8-schwartz}

\begin{definicja}[Przestrzeń Schwartza]
\label{def:pde8-schwartz}
Dla wielowskaźników $\alpha,\beta\in\N_0^d$ piszemy
$x^\alpha=\prod_{j=1}^d x_j^{\alpha_j}$ oraz
$\partial^\beta=\partial_1^{\beta_1}\cdots\partial_d^{\beta_d}$,
przyjmując $x_j^0=1$ także dla $x_j=0$.
Przestrzenią Schwartza na $\R^d$ nazywamy
\[
 \mathcal S(\R^d)=\bigl\{f\in C^\infty(\R^d;\C):
 p_{\alpha,\beta}(f)<\infty
 \text{ dla wszystkich }\alpha,\beta\in\N_0^d\bigr\},
\]
gdzie
\[
 p_{\alpha,\beta}(f)=\sup_{x\in\R^d}
 |x^\alpha\partial^\beta f(x)|.
\]
Wyposażamy ją w topologię wyznaczoną przez przeliczalną
rodzinę półnorm $p_{\alpha,\beta}$. Zbieżność $f_n\to f$
w $\mathcal S$ oznacza zbieżność do zera wszystkich
półnorm $p_{\alpha,\beta}(f_n-f)$.
\end{definicja}

Jest to zespolona przestrzeń wektorowa, ponieważ
liniowość pochodnych i nierówność trójkąta dają
$p_{\alpha,\beta}(af+bg)\le
|a|p_{\alpha,\beta}(f)+|b|p_{\alpha,\beta}(g)$.
Ta sama nierówność i jednorodność uzasadniają
własności półnorm. Rodzina rozdziela punkty,
gdyż $p_{0,0}(f)=0$ oznacza $f=0$.

\begin{lemat}[Równoważne półnormy i operacje]
\label{lem:pde8-seminorms}
Dla $N,m\in\N_0$ połóżmy
\[
 q_{N,m}(f)=\max_{|\beta|\le m}
 \sup_{x\in\R^d}(1+|x|)^N|\partial^\beta f(x)|.
\]
Rodzina $q_{N,m}$ określa tę samą topologię co
$p_{\alpha,\beta}$. Mnożenie przez $x_j$ oraz
różniczkowanie $\partial_j$ są ciągłymi operatorami
liniowymi $\mathcal S\to\mathcal S$.
Ponadto $\mathcal S(\R^d)\subset L^1(\R^d)\cap L^2(\R^d)$,
a oba włożenia są ciągłe.
\end{lemat}
\begin{proof}
Ponieważ $|x^\alpha|\le(1+|x|)^{|\alpha|}$, mamy
$p_{\alpha,\beta}(f)\le q_{|\alpha|,|\beta|}(f)$.
Odwrotnie, z $|x|\le\sum_{j=1}^d|x_j|$ i wzoru
wielomianowego dostajemy
\[
 (1+|x|)^N\le
 \sum_{|\alpha|\le N}
 \frac{N!}{(N-|\alpha|)!\alpha!}|x^\alpha|,
 \qquad \alpha!=\prod_{j=1}^d\alpha_j!.
\]
Mnożąc przez $|\partial^\beta f(x)|$, biorąc
supremum względem $x$ i maksimum po $|\beta|\le m$,
ograniczamy $q_{N,m}$ przez skończoną sumę półnorm
$p_{\alpha,\beta}$ z podanymi współczynnikami.
Obie rodziny wyznaczają więc tę samą topologię.

Z reguły Leibniza, dla wektora jednostkowego
$e_j$, mamy
\[
 \partial^\beta(x_j f)
 =x_j\partial^\beta f+\beta_j\partial^{\beta-e_j}f,
\]
gdzie drugi składnik pomijamy dla $\beta_j=0$.
Stąd
\[
 \begin{aligned}
 q_{N,m}(x_j f)&\le q_{N+1,m}(f)+m q_{N,m}(f),\\
 q_{N,m}(\partial_j f)&\le q_{N,m+1}(f).
 \end{aligned}
\]
Skończoność prawych stron zapewnia, że obrazy
należą do $\mathcal S$, a oszacowania dla każdej
półnormy dają ciągłość.

Dla $p=1,2$ wybieramy $N$ tak, aby $Np>d$.
Wtedy
\[
 \|f\|_{L^p}\le C_{d,N,p}q_{N,0}(f),
 \qquad
 C_{d,N,p}^p=\int_{\R^d}(1+|x|)^{-Np}
 \,\mathrm d\lambda_d(x)<\infty.
\]
Istotnie, na kuli jednostkowej całka jest skończona,
a poza nią całkowanie radialne sprowadza oszacowanie
do stałej razy $\int_1^\infty r^{d-1-Np}\,
\mathrm d\lambda_1(r)$, która jest skończona dla $Np>d$.
W wymiarze jeden całkujemy osobno na obu półprostych.
Daje to oba włożenia i ich ciągłość.
\end{proof}

Równoważne półnormy wyjaśniają szybki zanik:
dla każdego $N,\beta$ mamy
\[
 (1+|x|)^N|\partial^\beta f(x)|
 \le\frac{q_{N+1,|\beta|}(f)}{1+|x|}
 \longrightarrow 0\qquad (|x|\to\infty).
\]
Funkcja $e^{-|x|^2}$ należy do $\mathcal S$:
każda jej pochodna ma postać
$P_\beta(x)e^{-|x|^2}$ z wielomianem $P_\beta$.
Indukcja zaczyna się od $P_0=1$, a różniczkowanie
daje $P_{\beta+e_j}=\partial_j P_\beta-2x_jP_\beta$.
Dla każdej potęgi $k\ge0$ funkcja
$r^k e^{-r^2}$ jest ograniczona na $[0,\infty)$
i dąży do zera w nieskończoności (wykładnik
dominuje potęgi, jak w konstrukcji funkcji
testowej z sekcji~\ref{sec:pde7-tests}).
Każda funkcja z $C_c^\infty(\R^d)$ również należy
do $\mathcal S$, bo jej pochodne znikają poza
wspólnym zwartym nośnikiem.
Natomiast $(1+|x|^2)^{-1}$ nie należy do $\mathcal S$
w żadnym wymiarze $d\ge1$: dla
$x=te_1$ mamy $t^3/(1+t^2)\to+\infty$
przy $t\to+\infty$, więc $p_{3e_1,0}$ jest nieskończona.

\begin{stwierdzenie}[Zupełność przestrzeni Schwartza]
Przestrzeń $\mathcal S(\R^d)$ jest przestrzenią Frécheta,
czyli zupełną metryzowalną przestrzenią lokalnie wypukłą.
\end{stwierdzenie}
\begin{proof}
Rosnąca rodzina $Q_n=q_{n,n}$, $n\in\N_0$, wyznacza
tę samą topologię, ponieważ każda $q_{N,m}$ jest
ograniczona przez $Q_{\max\{N,m\}}$.
Metryka niezmiennicza względem przesunięć
\[
 D(f,g)=\sum_{n=0}^{\infty}2^{-n-1}
 \min\{1,Q_n(f-g)\}
\]
wyznacza tę topologię. Istotnie, nierówność
trójkąta wynika z półnorm i podaddytywności
$\min\{1,t\}$, a $Q_0$ rozdziela punkty.
Dla $0<\delta<1$ warunek
$D(f,g)<2^{-n-1}\delta$ daje $Q_n(f-g)<\delta$.
Odwrotnie, wybieramy $M$ z ogonem szeregu wag
mniejszym niż $\delta/2$; gdy $Q_M(f-g)<\delta/2$,
pierwsze $M+1$ składników mają sumę mniejszą
niż $\delta/2$, więc $D(f,g)<\delta$.
Te oszacowania charakteryzują również ciągi Cauchy'ego.

Niech $(f_\ell)$ będzie ciągiem Cauchy'ego w tej
metryce. Dla każdego $\beta$ ciąg
$\partial^\beta f_\ell$ spełnia jednostajny
warunek Cauchy'ego na $\R^d$ dzięki $Q_{|\beta|}$.
Ma zatem jednostajną granicę $g_\beta$, ciągłą
na $\R^d$. Dla każdego $x\in\R^d$, $t\in\R$
i $j\in\{1,\ldots,d\}$ mamy
\[
 \partial^\beta f_\ell(x+te_j)-\partial^\beta f_\ell(x)
 =\int_0^t\partial^{\beta+e_j}f_\ell(x+se_j)
 \,\mathrm d\lambda_1(s).
\]
Przechodzimy do granicy dzięki jednostajnej zbieżności
na odcinku (dla $t<0$ całka jest zorientowana).
Ciągłość $g_{\beta+e_j}$ daje
$\partial_j g_\beta=g_{\beta+e_j}$.
Indukcyjnie $f:=g_0$ jest gładka i
$\partial^\beta f=g_\beta$.

Ustalmy $n$ i $\varepsilon>0$. Dla
$\ell,k$ od pewnego indeksu mamy
$Q_n(f_\ell-f_k)<\varepsilon$. Dla każdego
$x$ i $|\beta|\le n$ przechodzimy z $k$
do nieskończoności w nierówności
\[
 (1+|x|)^n|\partial^\beta f_\ell(x)
              -\partial^\beta f_k(x)|<\varepsilon.
\]
Biorąc supremum i maksimum po przejściu do
granicy, otrzymujemy $Q_n(f_\ell-f)\le\varepsilon$.
W szczególności $Q_n(f)\le Q_n(f_\ell)+\varepsilon$
jest skończona. Ponieważ $n$ było dowolne,
$f\in\mathcal S$ i $f_\ell\to f$ we wszystkich
półnormach, a więc i w metryce $D$.
\end{proof}

\begin{twierdzenie}[Gęstość funkcji o zwartym nośniku]
\label{thm:pde8-density}
Przestrzeń $C_c^\infty(\R^d)$ jest gęsta w $\mathcal S(\R^d)$.
Jest też gęsta w $L^p(\R^d)$ dla $1\le p<\infty$.
\end{twierdzenie}
\begin{proof}
Wybierzmy funkcję odcinającą
$\psi\in C_c^\infty(\R^d)$ z $0\le\psi\le1$,
$\psi=1$ na $B_1(0)$ i $\psi=0$ poza $B_2(0)$.
Można przyjąć $\psi=\chi_{B_{3/2}(0)}*\rho_{1/4}$,
gdzie $\rho_\varepsilon$ jest przybliżoną jedynką
z sekcji~\ref{sec:pde7-mollifiers}.
Gładkość wynika z twierdzenia~\ref{thm:pde7-mollify}.
Nieujemność jądra i jego całka $1$ dają
$0\le\psi\le1$; dla $|x|\le1$ i $|y|\le1/4$
mamy $|x-y|<3/2$, a dla $|x|\ge2$ mamy
$|x-y|>3/2$. Stąd żądane wartości i zwarty nośnik.

Dla $f\in\mathcal S$ i $R\ge1$ określamy
$f_R(x)=\psi(x/R)f(x)\in C_c^\infty(\R^d)$.
Niech $\psi_R(x)=\psi(x/R)$.
Wszystkie pochodne $1-\psi_R$, włącznie z rzędem
zero, znikają dla $|x|<R$. Dla $|\gamma|>0$
ich moduły są ograniczone przez
$R^{-|\gamma|}\|\partial^\gamma\psi\|_\infty$;
dla $\gamma=0$ przez $1$.
Stosujemy wielowskaźnikową regułę Leibniza:
\[
 \partial^\beta(f-f_R)
 =\sum_{\gamma\le\beta}\binom{\beta}{\gamma}
 \partial^\gamma(1-\psi_R)\partial^{\beta-\gamma}f,
\]
gdzie $\gamma\le\beta$ oznacza nierówności
we wszystkich współrzędnych, a
$\binom{\beta}{\gamma}=\prod_j\binom{\beta_j}{\gamma_j}$.
Dla $|x|\ge R$ zachodzi
$(1+|x|)^N\le (1+R)^{-1}(1+|x|)^{N+1}$.
Zatem dla $|\beta|\le m$ każdy składnik po
pomnożeniu przez $(1+|x|)^N$ jest ograniczony
przez odpowiednią stałą zależną od $\psi,\beta,\gamma$
razy $(1+R)^{-1}q_{N+1,m}(f)$.
Sumując skończenie wiele składników, dostajemy
\[
 q_{N,m}(f-f_R)\le
 \frac{C_{\psi,m}}{1+R}q_{N+1,m}(f)
 \longrightarrow0.
\]
To daje gęstość w $\mathcal S$.

Dla $u\in L^p(\R^d)$, $1\le p<\infty$,
niech $v_R=u\chi_{B_R(0)}$. Z całkowalności
$|u|^p$ wynika
\[
 \|u-v_R\|_{L^p}^p
 =\int_{\R^d\setminus B_R(0)}|u(x)|^p
 \,\mathrm d\lambda_d(x)\longrightarrow0.
\]
Dla ustalonego $R$ funkcja
$v_R*\rho_\varepsilon$ jest gładka na $\R^d$
i dla $0<\varepsilon<1$ znika poza
$\overline{B_{R+1}(0)}$: jeśli $|x|>R+1$,
to dla $|y|\le R$ mamy $|x-y|>1>\varepsilon$.
Z twierdzenia~\ref{thm:pde7-approx-id}
na zwartym zbiorze $\overline{B_{R+1}(0)}$
wynika
\[
 \|v_R*\rho_\varepsilon-v_R\|_{L^p(\R^d)}
 \longrightarrow0\qquad(\varepsilon\downarrow0),
\]
ponieważ obie funkcje znikają poza tym zbiorem.
Dla dowolnego $\delta>0$ wybieramy najpierw
$R$ z $\|u-v_R\|_p<\delta/2$, następnie
$\varepsilon$ z błędem wygładzania mniejszym
niż $\delta/2$. Nierówność trójkąta daje
aproksymację $u$ z błędem mniejszym niż $\delta$
przez funkcję gładką o zwartym nośniku.
\end{proof}

\section{Transformata Fouriera na przestrzeni Schwartza}
\label{sec:pde8-transform}

\begin{twierdzenie}[Różniczkowanie i mnożenie]
\label{thm:pde8-rules}
Jeżeli $f\in\mathcal S(\R^d)$, to $\widehat f\in\mathcal S(\R^d)$
oraz dla każdego $\alpha\in\N_0^d$
\begin{equation}\label{eq:pde8-rules}
 \mathcal F(\partial^\alpha f)(\xi)=(i\xi)^\alpha\widehat f(\xi),
 \qquad \mathcal F(x^\alpha f)(\xi)=i^{|\alpha|}
 \partial_\xi^\alpha\widehat f(\xi).
\end{equation}
Ponadto $\mathcal F:\mathcal S(\R^d)\to\mathcal S(\R^d)$
jest ciągłym operatorem liniowym.
\end{twierdzenie}
\begin{proof}
Z lematu~\ref{lem:pde8-seminorms} wynika, że
$x^\beta f$ i $\partial^\beta f$ należą do
$\mathcal S(\R^d)$, a więc do $L^1(\R^d)$,
dla każdego wielowskaźnika $\beta$.
W szczególności całka definiująca $\widehat f$
jest bezwzględnie zbieżna.

Uzasadnijmy najpierw różniczkowanie względem $\xi$.
Dla $t\in\R$ mamy
\[
 |e^{-it}-1|
 =\left|\int_0^t(-i)e^{-is}\,\mathrm d\lambda_1(s)\right|
 \le |t|.
\]
Dla $h\ne0$ iloraz różnicowy
\[
 \frac{\widehat f(\xi+he_j)-\widehat f(\xi)}h
 =\int_{\R^d}e^{-ix\cdot\xi}
 \frac{e^{-ihx_j}-1}{h}f(x)\,\mathrm d\lambda_d(x)
\]
ma funkcję podcałkową ograniczoną co do modułu
przez całkowalną funkcję $|x_j f(x)|$.
Z twierdzenia o zbieżności zmajoryzowanej
otrzymujemy
\[
 \partial_{\xi_j}\widehat f(\xi)
 =\int_{\R^d}(-ix_j)e^{-ix\cdot\xi}f(x)
 \,\mathrm d\lambda_d(x).
\]
Ponieważ $x^\beta f\in L^1$ dla każdego $\beta$,
ten sam argument można zastosować do kolejnych
pochodnych. Indukcyjnie
\[
 \partial_\xi^\beta\widehat f(\xi)
 =(-i)^{|\beta|}\mathcal F(x^\beta f)(\xi).
\]
Każda z tych pochodnych jest ciągła: dla
$\xi_n\to\xi$ stosujemy twierdzenie o zbieżności
zmajoryzowanej z majorantą $|x^\beta f(x)|$.
Zatem $\widehat f\in C^\infty(\R^d;\C)$.
Mnożąc powyższą tożsamość przez $i^{|\beta|}$,
uzyskujemy drugą równość w~\eqref{eq:pde8-rules}.

Dla pierwszej równości ustalamy $j$ i pozostałe
współrzędne $x$. Całkowanie przez części względem
$x_j$ na przedziale $[-R,R]$ daje
\[
 \begin{aligned}
 \int_{-R}^R e^{-ix\cdot\xi}\partial_j f(x)
 \,\mathrm d\lambda_1(x_j)
 &=\left[e^{-ix\cdot\xi}f(x)\right]_{x_j=-R}^{x_j=R}\\
 &\quad+i\xi_j\int_{-R}^R e^{-ix\cdot\xi}f(x)
 \,\mathrm d\lambda_1(x_j).
 \end{aligned}
\]
Dla ustalonych pozostałych współrzędnych wyraz
brzegowy dąży do zera, ponieważ
$|f(x)|\le q_{1,0}(f)/(1+|x|)$.
Całki po prostej są bezwzględnie zbieżne dzięki
oszacowaniom $|f(x)|,|\partial_j f(x)|
\le q_{2,1}(f)(1+|x_j|)^{-2}$.
Po przejściu do granicy $R\to\infty$ całkujemy
względem pozostałych współrzędnych, stosując
twierdzenie Fubiniego: $f,\partial_j f\in L^1(\R^d)$.
Dla $d=1$ nie ma dodatkowego całkowania.
Otrzymujemy $\mathcal F(\partial_j f)=i\xi_j\widehat f$.
Ponieważ każda pochodna $f$ należy do $\mathcal S$,
iteracja daje pierwszą równość w~\eqref{eq:pde8-rules}.

Pozostaje wykazać skończoność półnorm transformaty
i jej ciągłość. Łącząc obie uzyskane tożsamości,
otrzymujemy dla dowolnych $\alpha,\beta\in\N_0^d$
\[
 \xi^\alpha\partial_\xi^\beta\widehat f(\xi)
 =i^{-|\alpha|}(-i)^{|\beta|}
 \mathcal F\bigl(\partial^\alpha(x^\beta f)\bigr)(\xi).
\]
Dla $v\in L^1(\R^d)$ zachodzi
$|\widehat v(\xi)|\le\|v\|_{L^1}$, ponieważ moduł
czynnika wykładniczego wynosi $1$.
Ustalmy całkowite $N>d$ i oznaczmy przez
$C_{d,N,1}$ stałą z lematu~\ref{lem:pde8-seminorms}.
Z wielowskaźnikowej reguły Leibniza wynika
\[
 \begin{aligned}
 p_{\alpha,\beta}(\widehat f)
 &\le \|\partial^\alpha(x^\beta f)\|_{L^1}\\
 &\le C_{d,N,1}
 \sum_{\substack{\gamma\le\alpha\\\gamma\le\beta}}
 \binom{\alpha}{\gamma}
 \frac{\beta!}{(\beta-\gamma)!}
 q_{N+|\beta|-|\gamma|,\,|\alpha-\gamma|}(f).
 \end{aligned}
\]
Tutaj pochodna $\partial^\gamma x^\beta$
jest zerowa, jeśli $\gamma\not\le\beta$;
w przeciwnym razie jest równa
$\beta!x^{\beta-\gamma}/(\beta-\gamma)!$.
Każda półnorma po prawej stronie jest
ograniczona przez $q_{N+|\beta|,|\alpha|}(f)$.
W konsekwencji istnieje skończona stała
$C_{\alpha,\beta,d,N}$, niezależna od $f$, taka że
\[
 p_{\alpha,\beta}(\widehat f)
 \le C_{\alpha,\beta,d,N}q_{N+|\beta|,|\alpha|}(f).
\]
Wszystkie półnormy $\widehat f$ są zatem skończone,
czyli $\widehat f\in\mathcal S$. Liniowość wynika
z liniowości całki, a uzyskane oszacowania
każdej półnormy dają ciągłość operatora.
\end{proof}

\begin{stwierdzenie}[Lemat Riemanna--Lebesgue'a]
\label{prop:pde8-RL}
Jeśli $f\in L^1(\R^d)$, to $\widehat f$ jest ograniczona,
jednostajnie ciągła i $\widehat f(\xi)\to0$
przy $|\xi|\to\infty$.
\end{stwierdzenie}
\begin{proof}
Całka definiująca $\widehat f$ nie zależy od wyboru
reprezentanta klasy $f$ i daje
$|\widehat f(\xi)|\le\|f\|_{L^1}$.
Dla $\xi,h\in\R^d$ mamy
\[
 |\widehat f(\xi+h)-\widehat f(\xi)|
 \le\int_{\R^d}|e^{-ix\cdot h}-1|\,|f(x)|
 \,\mathrm d\lambda_d(x).
\]
Prawa strona nie zależy od $\xi$ i dąży do zera
dla $h\to0$: funkcja podcałkowa dąży punktowo
do zera i jest ograniczona przez $2|f(x)|$.
Twierdzenie o zbieżności zmajoryzowanej daje
więc jednostajną ciągłość.

Ustalmy $\varepsilon>0$. Z twierdzenia
~\ref{thm:pde8-density} wybieramy
$g\in C_c^\infty(\R^d)$ z
$\|f-g\|_{L^1}<\varepsilon/2$.
Ponieważ $g\in\mathcal S(\R^d)$, poprzednie
twierdzenie daje $\widehat g\in\mathcal S(\R^d)$.
Istnieje zatem $R>0$ takie, że
$|\widehat g(\xi)|<\varepsilon/2$ dla $|\xi|>R$.
Dla tych $\xi$ z liniowości i oszacowania
transformaty funkcji całkowalnej wynika
\[
 |\widehat f(\xi)|
 \le\|f-g\|_{L^1}+|\widehat g(\xi)|<\varepsilon.
\]
Dowodzi to zaniku w nieskończoności.
\end{proof}

\section{Transformata funkcji Gaussa i wzór odwrócenia}
\label{sec:pde8-inversion}

\begin{lemat}[Transformata funkcji Gaussa]
\label{lem:pde8-gauss}
Dla $a>0$ i $\xi\in\R^d$ zachodzi
\begin{equation}\label{eq:pde8-gauss}
 \mathcal F(e^{-a|x|^2})(\xi)
 =\left(\frac\pi a\right)^{d/2}e^{-|\xi|^2/(4a)}.
\end{equation}
\end{lemat}
\begin{proof}
W wymiarze jeden niech $g_a(x)=e^{-ax^2}$.
Funkcja ta należy do $\mathcal S(\R)$:
argument z sekcji~\ref{sec:pde8-schwartz} stosuje się
z rekurencją wielomianów
$P_{k+1}(x)=P_k'(x)-2axP_k(x)$.
Niech $I=\widehat{g_a}$. Ponieważ
$g_a'=-2axg_a$, tożsamości~\eqref{eq:pde8-rules}
dają
\[
 i\xi I(\xi)=\widehat{g_a'}(\xi)
 =-2a\mathcal F(xg_a)(\xi)=-2aiI'(\xi).
\]
Stąd $I'(\xi)=-\xi I(\xi)/(2a)$ i
\[
 \frac{\mathrm d}{\mathrm d\xi}
 \left(e^{\xi^2/(4a)}I(\xi)\right)=0.
\]
Zatem $I(\xi)=I(0)e^{-\xi^2/(4a)}$.

Dla wyznaczenia $I(0)$ oznaczmy
$J=\int_\R e^{-t^2}\,\mathrm d\lambda_1(t)$.
Jest to liczba dodatnia i skończona, ponieważ
dla $|t|\ge1$ mamy $e^{-t^2}\le e^{-|t|}$.
Z twierdzenia Tonellego oraz całkowania w
zmiennych biegunowych na $\R^2$ otrzymujemy
\[
 J^2=\int_{\R^2}e^{-(x^2+y^2)}\,\mathrm d\lambda_2(x,y)
 =2\pi\int_0^\infty e^{-r^2}r\,\mathrm d\lambda_1(r)
 =\pi.
\]
Ponieważ $J>0$, mamy $J=\sqrt\pi$.
Podstawienie $t=\sqrt a\,x$ daje
$I(0)=a^{-1/2}J=\sqrt{\pi/a}$.
Nie korzystaliśmy tu ze wzoru odwrócenia.

W wymiarze $d$ funkcja podcałkowa jest
iloczynem czynników jednowymiarowych, a całka
z jej modułu jest skończona i wynosi
$(\pi/a)^{d/2}$ (Tonelli). Z twierdzenia
Fubiniego wynika
\[
 \begin{aligned}
 \mathcal F(e^{-a|x|^2})(\xi)
 &=\prod_{j=1}^d
 \int_\R e^{-ax_j^2}e^{-ix_j\xi_j}\,
 \mathrm d\lambda_1(x_j)\\
 &=\left(\frac{\pi}{a}\right)^{d/2}
 e^{-|\xi|^2/(4a)}.
 \end{aligned}
\]
\end{proof}

\begin{twierdzenie}[Odwrócenie transformaty]
\label{thm:pde8-inversion}
Dla $f\in\mathcal S(\R^d)$ zachodzi
\[
 f(x)=(2\pi)^{-d}\int_{\R^d}e^{ix\cdot\xi}
 \widehat f(\xi)\,\mathrm d\lambda_d(\xi),
 \qquad x\in\R^d.
\]
Zatem $\mathcal F$ jest topologicznym automorfizmem
$\mathcal S(\R^d)$, którego odwrotność podano w
\eqref{eq:pde8-convention}.
\end{twierdzenie}
\begin{proof}
Dla $\varepsilon>0$ określmy
\[
 I_\varepsilon(x)=(2\pi)^{-d}\int_{\R^d}
 e^{ix\cdot\xi}e^{-\varepsilon|\xi|^2/2}
 \widehat f(\xi)\,\mathrm d\lambda_d(\xi).
\]
Po wstawieniu definicji $\widehat f$ całka
podwójna jest bezwzględnie zbieżna, ponieważ
całka z modułu (bez stałej $(2\pi)^{-d}$) wynosi
\[
 \|f\|_{L^1}
 \int_{\R^d}e^{-\varepsilon|\xi|^2/2}
 \,\mathrm d\lambda_d(\xi)
 =\|f\|_{L^1}(2\pi/\varepsilon)^{d/2}<\infty.
\]
Stosując Fubiniego i lemat~\ref{lem:pde8-gauss}
z $a=\varepsilon/2$, otrzymujemy
\[
 \begin{aligned}
 I_\varepsilon(x)
 &=\int_{\R^d}f(y)G_\varepsilon(x-y)
 \,\mathrm d\lambda_d(y),\\
 G_\varepsilon(z)
 &:=(2\pi)^{-d}\int_{\R^d}e^{iz\cdot\xi}
 e^{-\varepsilon|\xi|^2/2}\,\mathrm d\lambda_d(\xi)\\
 &=(2\pi\varepsilon)^{-d/2}
 e^{-|z|^2/(2\varepsilon)}.
 \end{aligned}
\]
Jądro jest nieujemne i ma całkę $1$, co wynika
z obliczenia całki Gaussa. Nie ma jednak zwartego
nośnika, dlatego uzasadnimy jego własność
przybliżonej jedynki bezpośrednio.

Po podstawieniu $y=x-\sqrt\varepsilon\,z$ mamy
\[
 I_\varepsilon(x)-f(x)
 =\int_{\R^d}G_1(z)
 \bigl(f(x-\sqrt\varepsilon\,z)-f(x)\bigr)
 \,\mathrm d\lambda_d(z).
\]
Niech $M=\sup_{x\in\R^d}|\nabla f(x)|<\infty$,
gdzie dla gradientu zespolonego stosujemy
normę euklidesową. Całkowanie wzdłuż odcinka
daje $|f(x-h)-f(x)|\le M|h|$ dla $h\in\R^d$.
W konsekwencji
\[
 \sup_x|I_\varepsilon(x)-f(x)|
 \le M\sqrt\varepsilon
 \int_{\R^d}|z|G_1(z)\,\mathrm d\lambda_d(z)
 \longrightarrow0.
\]
Całka z pierwszego momentu jest skończona
dzięki wykładniczemu zanikowi Gaussa
(w zmiennych biegunowych wystarcza całkowalność
$r^d e^{-r^2/2}$ na $(0,\infty)$; dla $d=1$
całkujemy na obu półprostych).

Z drugiej strony $\widehat f\in L^1$ na mocy
twierdzenia~\ref{thm:pde8-rules} i
lematu~\ref{lem:pde8-seminorms}. Mamy
\[
 \begin{aligned}
 &\sup_x\left| I_\varepsilon(x)
 -(2\pi)^{-d}\int_{\R^d}e^{ix\cdot\xi}
 \widehat f(\xi)\,\mathrm d\lambda_d(\xi)\right|\\
 &\quad\le (2\pi)^{-d}\int_{\R^d}
 (1-e^{-\varepsilon|\xi|^2/2})|\widehat f(\xi)|
 \,\mathrm d\lambda_d(\xi)\longrightarrow0
 \end{aligned}
\]
z twierdzenia o zbieżności zmajoryzowanej
z majorantą $(2\pi)^{-d}|\widehat f(\xi)|$.
Obie granice są więc równe, co dowodzi wzoru.

Sprawdźmy jeszcze odwracalność w obu kierunkach.
Niech $\mathcal R f(x)=f(-x)$. Operator
$\mathcal R$ zachowuje każdą półnormę
$p_{\alpha,\beta}$, spełnia $\mathcal R^2=I$
i komutuje z $\mathcal F$:
podstawienie $y=-x$ daje
$\mathcal F(\mathcal R f)(\xi)=\widehat f(-\xi)$.
Operator określony całką odwrotną to
$\mathcal G=(2\pi)^{-d}\mathcal R\mathcal F$,
więc odwzorowuje $\mathcal S$ ciągle w siebie.
Udowodniony wzór mówi $\mathcal G\mathcal F=I$,
a komutowanie $\mathcal R$ i $\mathcal F$ daje
$\mathcal F\mathcal G=\mathcal G\mathcal F=I$.
Zatem $\mathcal G=\mathcal F^{-1}$ i oba
operatory są ciągłe. Otrzymujemy ponadto
$\mathcal F^2=(2\pi)^d\mathcal R$.
\end{proof}

\begin{wniosek}[Wzory na splot i iloczyn]
\label{cor:pde8-convolution}
Dla $f,g\in\mathcal S(\R^d)$ funkcje $f*g$ oraz
$fg$ należą do $\mathcal S(\R^d)$ i
\[
 \widehat{f*g}=\widehat f\,\widehat g,\qquad
 \widehat{fg}=(2\pi)^{-d}(\widehat f*\widehat g).
\]
\end{wniosek}
\begin{proof}
Przyjmujemy $(f*g)(x)=\int_{\R^d}f(y)g(x-y)
\,\mathrm d\lambda_d(y)$. Całka jest bezwzględnie
zbieżna dla każdego $x$, ponieważ $f\in L^1$,
a $g$ jest ograniczona.
Różniczkowanie pod całką daje
$\partial^\beta(f*g)=f*(\partial^\beta g)$:
ilorazy różnicowe są zmajoryzowane przez
$|f(y)|$ razy normę supremum odpowiedniej
pochodnej $g$ wyższego rzędu, a ciągłość każdej
pochodnej wynika z tej samej metody z majorantą
$|f(y)|\|\partial^\beta g\|_\infty$.
Z nierówności
$1+|x|\le(1+|y|)(1+|x-y|)$ wynika
\[
 q_{N,m}(f*g)
 \le q_{N,m}(g)
 \int_{\R^d}(1+|y|)^N|f(y)|\,\mathrm d\lambda_d(y).
\]
Ostatnia całka jest skończona, gdyż dla
całkowitego $L>d$ funkcja podcałkowa jest
ograniczona przez $q_{N+L,0}(f)(1+|y|)^{-L}$.
Zatem $f*g\in\mathcal S$.
Z reguły Leibniza, rozdzielając wagę
$(1+|x|)^N$ na pierwszy czynnik,
otrzymujemy także
$q_{N,m}(fg)\le 2^m q_{N,m}(f)q_{0,m}(g)$,
więc $fg\in\mathcal S$.

Całka podwójna z modułu $f(y)g(x-y)$
wynosi $\|f\|_{L^1}\|g\|_{L^1}$ dzięki
niezmienniczości miary względem przesunięć
i twierdzeniu Tonellego. Zatem Fubini
i podstawienie $z=x-y$ dają
\[
 \begin{aligned}
 \widehat{f*g}(\xi)
 &=\int_{\R^d}f(y)
 \int_{\R^d}e^{-ix\cdot\xi}g(x-y)
 \,\mathrm d\lambda_d(x)\,\mathrm d\lambda_d(y)\\
 &=\int_{\R^d}e^{-iy\cdot\xi}f(y)
 \widehat g(\xi)\,\mathrm d\lambda_d(y)
 =\widehat f(\xi)\widehat g(\xi).
 \end{aligned}
\]
Dla drugiej tożsamości podstawiamy wzór
odwrócenia dla $g$. Bezwzględną całkowalność
całki podwójnej zapewnia
$\|f\|_{L^1}\|\widehat g\|_{L^1}<\infty$.
Stąd
\[
 \begin{aligned}
 \widehat{fg}(\xi)
 &=(2\pi)^{-d}\int_{\R^d}\widehat g(\eta)
 \int_{\R^d}e^{-ix\cdot(\xi-\eta)}f(x)
 \,\mathrm d\lambda_d(x)\,\mathrm d\lambda_d(\eta)\\
 &=(2\pi)^{-d}\int_{\R^d}
 \widehat f(\xi-\eta)\widehat g(\eta)
 \,\mathrm d\lambda_d(\eta)
 =(2\pi)^{-d}(\widehat f*\widehat g)(\xi).
 \end{aligned}
\]
W przyjętej konwencji czynnik $(2\pi)^{-d}$
występuje we wzorze na transformatę iloczynu,
a nie splotu.
\end{proof}

\section[Twierdzenie Plancherela i przestrzeń L2]{Twierdzenie Plancherela i przestrzeń $L^2$}
\label{sec:pde8-plancherel}

\begin{twierdzenie}[Tożsamość Parsevala--Plancherela]
\label{thm:pde8-plancherel}
Dla $f,g\in\mathcal S(\R^d)$ zachodzi
\begin{equation}\label{eq:pde8-parseval}
 \int_{\R^d}f(x)\overline{g(x)}\,\mathrm d\lambda_d(x)
 =(2\pi)^{-d}\int_{\R^d}\widehat f(\xi)
 \overline{\widehat g(\xi)}\,\mathrm d\lambda_d(\xi).
\end{equation}
Transformata Fouriera przedłuża się jednoznacznie do
ciągłego operatora liniowego
$\mathcal F:L^2(\R^d)\to L^2(\R^d)$, dla którego
$\|\widehat f\|_2=(2\pi)^{d/2}\|f\|_2$.
Operator $(2\pi)^{-d/2}\mathcal F$ jest unitarny,
a tożsamość~\eqref{eq:pde8-parseval} zachodzi także
dla dowolnych $f,g\in L^2(\R^d)$.
\end{twierdzenie}
\begin{proof}
Ze wzoru odwrócenia z twierdzenia
~\ref{thm:pde8-inversion} mamy
\[
 \overline{g(x)}=(2\pi)^{-d}\int_{\R^d}
 e^{-ix\cdot\xi}\overline{\widehat g(\xi)}
 \,\mathrm d\lambda_d(\xi).
\]
Po pomnożeniu przez $f(x)$ całka podwójna
jest bezwzględnie zbieżna, ponieważ całka
z modułu, bez czynnika $(2\pi)^{-d}$, wynosi
$\|f\|_1\|\widehat g\|_1<\infty$.
Stąd z twierdzenia Fubiniego
\[
 \begin{aligned}
 \int_{\R^d}f(x)\overline{g(x)}\,\mathrm d\lambda_d(x)
 &=(2\pi)^{-d}\int_{\R^d}
 \overline{\widehat g(\xi)}
 \int_{\R^d}e^{-ix\cdot\xi}f(x)
 \,\mathrm d\lambda_d(x)\,\mathrm d\lambda_d(\xi)\\
 &=(2\pi)^{-d}\int_{\R^d}
 \widehat f(\xi)\overline{\widehat g(\xi)}
 \,\mathrm d\lambda_d(\xi).
 \end{aligned}
\]
W szczególności dla $f=g$ otrzymujemy
$\|\widehat f\|_2^2=(2\pi)^d\|f\|_2^2$.
W zespolonej przestrzeni Hilberta $L^2(\R^d)$
przyjmujemy iloczyn skalarny liniowy względem
pierwszego argumentu,
$\langle f,g\rangle_2=\int_{\R^d}f(x)\overline{g(x)}
\,\mathrm d\lambda_d(x)$.

Oznaczmy przez $U_0=(2\pi)^{-d/2}\mathcal F$
operator na $\mathcal S$. Jest on liniową izometrią
w normie $L^2$. Dla dowolnego $f\in L^2$
wybieramy ciąg $f_n\in\mathcal S$ z
$\|f_n-f\|_2\to0$, korzystając z twierdzenia
~\ref{thm:pde8-density}. Wówczas
\[
 \|U_0f_n-U_0f_k\|_2=\|f_n-f_k\|_2.
\]
Ciąg $(U_0f_n)$ jest ciągiem Cauchy'ego,
więc ma granicę w zupełnej przestrzeni $L^2$.
Definiujemy $Uf=\lim_n U_0f_n$ w $L^2$.
Jeśli $g_n\in\mathcal S$ jest innym ciągiem
aproksymującym $f$, to
$\|U_0f_n-U_0g_n\|_2=\|f_n-g_n\|_2
\le\|f_n-f\|_2+\|f-g_n\|_2\to0$.
Definicja nie zależy więc od wyboru ciągu.

Liniowość $U$ wynika z liniowości $U_0$:
dla kombinacji liniowej można użyć odpowiedniej
kombinacji ciągów aproksymujących. Ciągłość normy
daje $\|Uf\|_2=\lim_n\|f_n\|_2=\|f\|_2$.
Operator $U$ przedłuża $U_0$ (dla $f\in\mathcal S$
można wziąć ciąg stały), a jednoznaczność
ciągłego przedłużenia wynika z gęstości $\mathcal S$.

Obraz $U(L^2)$ jest domknięty: jeśli $Uf_n\to h$
w $L^2$, to $(f_n)$ jest ciągiem Cauchy'ego
dzięki izometryczności $U$. Istnieje zatem
$f\in L^2$ z $f_n\to f$, a ciągłość $U$ daje
$h=Uf$. Zarazem $U(\mathcal S)=\mathcal S$,
gdyż $\mathcal F$ jest automorfizmem $\mathcal S$
i mnożenie przez niezerową stałą nie zmienia
tej przestrzeni. Obraz $U(L^2)$ zawiera więc
gęsty podzbiór $\mathcal S$, a będąc domknięty,
jest całym $L^2$.

Dla $f,g\in L^2$ wybierzmy $f_n,g_n\in\mathcal S$
zbiegające do nich w normie $L^2$.
Z nierówności Cauchy'ego--Schwarza wynika
\[
 |\langle f_n,g_n\rangle_2-\langle f,g\rangle_2|
 \le\|f_n-f\|_2\|g_n\|_2+\|f\|_2\|g_n-g\|_2
 \longrightarrow0,
\]
ponieważ ciąg $(\|g_n\|_2)$ jest ograniczony.
Analogiczne oszacowanie dla $Uf_n,Ug_n$
uzasadnia przejście do granicy w tożsamości
$\langle U_0f_n,U_0g_n\rangle_2
=\langle f_n,g_n\rangle_2$.
Zatem $U$ zachowuje iloczyn skalarny i jest
surjektywny, czyli jest operatorem unitarnym.
Przyjmując $\mathcal F=(2\pi)^{d/2}U$ na $L^2$,
otrzymujemy wszystkie twierdzone własności.
\end{proof}

Operator odwrotny na $L^2$ ma postać
\[
 \mathcal F^{-1}=(2\pi)^{-d}\mathcal R\mathcal F,
 \qquad \mathcal R f(x)=f(-x).
\]
Rzeczywiście, $\mathcal R$ jest izometrią $L^2$
na mocy podstawienia $y=-x$. Obie strony
tożsamości są ciągłe na $L^2$ i pokrywają się
na gęstej przestrzeni $\mathcal S$ dzięki
twierdzeniu~\ref{thm:pde8-inversion}.
W szczególności także na $L^2$ mamy
$\mathcal F^2=(2\pi)^d\mathcal R$.

\begin{stwierdzenie}[Zgodność transformaty całkowej i przedłużenia]
Jeśli $f\in L^1(\R^d)\cap L^2(\R^d)$, to
transformata określona całką
\eqref{eq:pde8-convention} jest reprezentantem
klasy $\mathcal Ff\in L^2$ z twierdzenia
Plancherela. Analogiczna zgodność zachodzi
dla transformaty odwrotnej.
\end{stwierdzenie}
\begin{proof}
Potrzebujemy aproksymacji jednocześnie w obu normach.
Dla $R>0$ połóżmy $v_R=f\chi_{B_R(0)}$.
Z całkowalności $|f|$ i $|f|^2$ wynika
$\|f-v_R\|_1+\|f-v_R\|_2\to0$.
Dla ustalonego $R$ wygładzamy $v_R$ przez
jądro $\rho_\varepsilon$ z sekcji
~\ref{sec:pde7-mollifiers}. Tak jak w dowodzie
twierdzenia~\ref{thm:pde8-density},
$v_R*\rho_\varepsilon\in C_c^\infty(\R^d)$,
a dla $0<\varepsilon<1$ obie funkcje
$v_R$ i $v_R*\rho_\varepsilon$ znikają
prawie wszędzie poza $\overline{B_{R+1}(0)}$.
Twierdzenie~\ref{thm:pde7-approx-id} dla
$p=1$ i $p=2$ daje zatem zbieżność w obu
normach globalnych. Wybierając najpierw $R_n$,
a następnie $\varepsilon_n$, uzyskujemy
$f_n\in C_c^\infty$ takie, że
$\|f-f_n\|_1+\|f-f_n\|_2<1/n$.

Oznaczmy przez $h$ transformatę całkową $f$,
a przez $H$ jej transformatę w sensie $L^2$.
Z oszacowania dla transformaty całkowej
mamy $\|\widehat f_n-h\|_\infty\le\|f_n-f\|_1\to0$.
Z twierdzenia Plancherela
$\|\widehat f_n-H\|_2=(2\pi)^{d/2}\|f_n-f\|_2\to0$.
Dla każdej kuli $B_R(0)$ otrzymujemy więc
\[
 \|H-h\|_{L^2(B_R(0))}
 \le\|H-\widehat f_n\|_2
 +\lambda_d(B_R(0))^{1/2}\|\widehat f_n-h\|_\infty
 \longrightarrow0.
\]
Funkcje $H,h$ są zatem równe prawie wszędzie
na każdej takiej kuli, a przeliczalne pokrycie
$\R^d$ kulami o całkowitych promieniach daje
$H=h$ prawie wszędzie na całej przestrzeni.
Dla odwrotności stosujemy wykazaną tożsamość
$\mathcal F^{-1}=(2\pi)^{-d}\mathcal R\mathcal F$
oraz tę samą tożsamość dla całek na $L^1$.
\end{proof}

\begin{uwaga}[Właściwe znaczenie całki dla $L^2$]
Jeśli $f\in L^2\setminus L^1$, zapis $\widehat f$
oznacza przedłużenie operatora z twierdzenia
Plancherela, a nie automatycznie zbieżną całkę
z~\eqref{eq:pde8-convention}.
Dla każdego $R>0$ funkcja
$f_R=f\chi_{B_R(0)}$ należy do $L^1\cap L^2$,
gdyż $\|f_R\|_1\le
\lambda_d(B_R(0))^{1/2}\|f\|_2$.
Ponadto $f_R\to f$ w $L^2$, więc
\[
 \widehat f
 =\lim_{R\to\infty}^{L^2}
 \int_{B_R(0)} e^{-ix\cdot\xi}f(x)
 \,\mathrm d\lambda_d(x).
\]
Nie jest to twierdzenie o zbieżności punktowej
tej rodziny całek.

Na przykład $f(x)=(1+|x|)^{-3/4}$ na $\R$ spełnia
\[
 \int_\R |f(x)|^2\,\mathrm d\lambda_1(x)
 =2\int_0^\infty(1+x)^{-3/2}\,\mathrm d\lambda_1(x)=4,
\]
lecz $\int_\R |f(x)|\,\mathrm d\lambda_1(x)=+\infty$,
bo całka $\int_0^\infty(1+x)^{-3/4}\,\mathrm d\lambda_1(x)$
jest rozbieżna. W szczególności całka Fouriera
dla $\xi=0$ jest rozbieżna, mimo że transformata
w sensie $L^2$ istnieje jako klasa funkcji.
\end{uwaga}

\section[Dystrybucje temperowane]{Dystrybucje temperowane: przestrzeń sprzężona do $\mathcal S$}
\label{sec:pde8-tempered}

\begin{definicja}[Dystrybucja temperowana]
\label{def:pde8-tempered}
Przestrzenią dystrybucji temperowanych $\mathcal S'(\R^d)$
nazywamy ciągłą przestrzeń sprzężoną do $\mathcal S(\R^d)$,
złożoną ze wszystkich ciągłych zespolonych funkcjonałów
liniowych $T:\mathcal S(\R^d)\to\C$.
Działanie funkcjonału zapisujemy jako
$\langle T,\varphi\rangle=T(\varphi)$.
Jest ono liniowe względem $\varphi$, bez sprzężenia
zespolonego, zgodnie z konwencją rozdziału
~\ref{chap:pde-distributions}.
Zbieżność $T_n\to T$ w sensie dystrybucji temperowanych
oznacza zbieżność
\[
 \langle T_n,\varphi\rangle\to\langle T,\varphi\rangle
 \qquad\text{dla każdego }\varphi\in\mathcal S(\R^d).
\]
Jest to zbieżność w topologii słabej*,
czyli $\sigma(\mathcal S',\mathcal S)$.
Nie oznacza ona zbieżności jednostajnej na wszystkich
funkcjach testowych.
\end{definicja}

\begin{twierdzenie}[Globalne oszacowanie dystrybucji temperowanej]
\label{thm:pde8-growth}
Funkcjonał liniowy $T$ na $\mathcal S(\R^d)$ jest ciągły
dokładnie wtedy, gdy istnieją $C>0$ i $N,m\in\N_0$ takie, że
\begin{equation}\label{eq:pde8-growth}
 |\langle T,\varphi\rangle|\le C\,q_{N,m}(\varphi)
 \qquad(\varphi\in\mathcal S(\R^d)).
\end{equation}
\end{twierdzenie}
\begin{proof}
Jeśli $T$ jest ciągły, to przeciwobraz zbioru
$\{z\in\C:|z|<1\}$ jest otoczeniem zera.
Z definicji topologii wyznaczonej przez półnormy
istnieje skończona rodzina warunków
$q_{N_j,m_j}(\varphi)<r_j$, gdzie $r_j>0$,
których jednoczesne spełnienie zapewnia
$|\langle T,\varphi\rangle|<1$.
Wybieramy $N=\max_j N_j$, $m=\max_j m_j$
i $r=\min_j r_j>0$.
Ponieważ $q_{N_j,m_j}\le q_{N,m}$, wystarcza
$q_{N,m}(\varphi)<r$.
Dla $\varphi\ne0$ mamy $q_{N,m}(\varphi)>0$,
gdyż $q_{N,m}(\varphi)\ge\|\varphi\|_\infty$.
Stosując powyższy warunek do
$r\varphi/(2q_{N,m}(\varphi))$, otrzymujemy
\[
 |\langle T,\varphi\rangle|
 <\frac{2}{r}q_{N,m}(\varphi).
\]
Dla $\varphi=0$ nierówność nieostra zachodzi
z liniowości. To daje~\eqref{eq:pde8-growth}.
Odwrotnie, jeśli zachodzi~\eqref{eq:pde8-growth},
to dla dowolnego $\eta>0$ warunek
$q_{N,m}(\varphi)<\eta/C$ pociąga
$|T(\varphi)|<\eta$. Funkcjonał jest więc ciągły
w zerze, a z liniowości ciągły w każdym punkcie.
\end{proof}

\begin{stwierdzenie}[Granice dystrybucji temperowanych]
Jeśli $T_n\in\mathcal S'(\R^d)$ i dla każdego
$\varphi\in\mathcal S(\R^d)$ istnieje skończona
granica $T(\varphi)=\lim_n T_n(\varphi)$,
to $T\in\mathcal S'(\R^d)$. Co więcej, istnieją
$C>0$ oraz $N,m\in\N_0$ wspólne dla całego ciągu,
takie że
\[
 |T_n(\varphi)|\le C q_{N,m}(\varphi)
 \qquad(n\ge1,\ \varphi\in\mathcal S(\R^d)).
\]
\end{stwierdzenie}
\begin{proof}
Dla $k\ge1$ określmy zbiory
\[
 A_k=\{\varphi\in\mathcal S:
 |T_n(\varphi)|\le k\text{ dla każdego }n\ge1\}.
\]
Są one domknięte jako przecięcia domkniętych
przeciwobrazów zbioru $\{z:|z|\le k\}$ przez
ciągłe funkcjonały $T_n$. Ich suma jest równa
$\mathcal S$, ponieważ każdy zbieżny ciąg
liczb zespolonych jest ograniczony.
Przestrzeń $\mathcal S$ jest zupełna
metryzowalna (sekcja~\ref{sec:pde8-schwartz}),
więc z twierdzenia Baire'a pewien $A_k$
ma niepuste wnętrze. Istnieją zatem
$\varphi_0\in A_k$, $r>0$ i $N,m\in\N_0$
takie, że $q_{N,m}(\psi)<r$ pociąga
$\varphi_0+\psi\in A_k$. Możemy użyć jednej
półnormy dzięki monotoniczności rodziny
$q_{N,m}$, tak jak w poprzednim dowodzie.
Wtedy dla każdego $n$
\[
 |T_n(\psi)|\le
 |T_n(\varphi_0+\psi)|+|T_n(\varphi_0)|\le2k.
\]
Dla $\varphi\ne0$ podstawiamy
$\psi=r\varphi/(2q_{N,m}(\varphi))$,
uzyskując $|T_n(\varphi)|\le(4k/r)q_{N,m}(\varphi)$.
Dla $\varphi=0$ oszacowanie jest oczywiste.
Przechodząc do granicy, dostajemy tę samą
nierówność dla $T$. Granica punktowa
funkcjonałów liniowych jest liniowa, więc
twierdzenie~\ref{thm:pde8-growth} daje ciągłość $T$.
\end{proof}

\begin{uwaga}[Relacja między $\mathcal D'$ a $\mathcal S'$]
Ograniczenie $T\in\mathcal S'(\R^d)$ do
$C_c^\infty(\R^d)$ jest dystrybucją.
Istotnie, jeśli $\supp\varphi\subset K$ dla
niepustego zbioru zwartego $K$, to
\[
 q_{N,m}(\varphi)
 \le (1+\max_{x\in K}|x|)^N
 \max_{|\beta|\le m}\sup_{x\in K}
 |\partial^\beta\varphi(x)|.
\]
Dla $K=\varnothing$ jedyną taką funkcją jest
funkcja zerowa. Globalne oszacowanie
~\eqref{eq:pde8-growth} daje zatem lokalne
oszacowanie z twierdzenia
~\ref{thm:pde7-local-order}.
Odwzorowanie ograniczenia
$\mathcal S'\to\mathcal D'$ jest liniowe
i iniektywne: jeśli dwa ciągłe funkcjonały
na $\mathcal S$ zgadzają się na $C_c^\infty$,
to zgadzają się na całej $\mathcal S$
dzięki gęstości z twierdzenia~\ref{thm:pde8-density}.
Możemy więc utożsamiać $\mathcal S'$
z podprzestrzenią $\mathcal D'$.
Dystrybucja jest temperowana dokładnie wtedy,
gdy posiada ciągłe przedłużenie z $C_c^\infty$
na $\mathcal S$; takie przedłużenie jest jednoznaczne.
Nie wszystkie dystrybucje są temperowane,
jak pokaże ostatni z poniższych przykładów.
W szczególności transformata Fouriera dowolnej
dystrybucji z $\mathcal D'$ nie ma automatycznie
sensu w rachunku, który zbudujemy dla $\mathcal S'$.
\end{uwaga}

\begin{przyklad}[Obiekty temperowane i nietemperowane]
\begin{enumerate}
\item Niech $f$ będzie funkcją mierzalną spełniającą
$|f(x)|\le C(1+|x|)^M$ prawie wszędzie,
gdzie $C>0$, $M\in\N_0$.
Definiujemy
\[
 \langle T_f,\varphi\rangle
 =\int_{\R^d}f(x)\varphi(x)\,\mathrm d\lambda_d(x).
\]
Dla całkowitego $L>d$ mamy
\[
 \int_{\R^d}|f(x)\varphi(x)|\,\mathrm d\lambda_d(x)
 \le C C_{d,L,1}\,q_{M+L,0}(\varphi),
\]
gdzie $C_{d,L,1}$ jest stałą z lematu
~\ref{lem:pde8-seminorms}.
Całka jest zatem bezwzględnie zbieżna,
a liniowy funkcjonał $T_f$ jest ciągły
z twierdzenia~\ref{thm:pde8-growth}.
Należą tu wielomiany, stałe i funkcje istotnie
ograniczone. Punktowe oszacowanie wielomianowe
jest warunkiem wystarczającym, nie definicją
temperowaności funkcji.

\item Każda funkcja $f\in L^p(\R^d)$,
$1\le p\le\infty$, również określa
dystrybucję temperowaną. Niech $p'$ będzie
wykładnikiem sprzężonym, czyli
$1/p+1/p'=1$ z konwencją $1/\infty=0$.
Z nierówności Höldera mamy
$|\langle T_f,\varphi\rangle|\le
\|f\|_p\|\varphi\|_{p'}$.
Jeśli $p'<\infty$, wybieramy całkowite $N$
takie, że $Np'>d$. Z tego samego całkowania
radialnego co w lemacie~\ref{lem:pde8-seminorms}
wynika $\|\varphi\|_{p'}\le
C_{d,N,p'}q_{N,0}(\varphi)$, gdzie
$C_{d,N,p'}^{p'}=\int_{\R^d}(1+|x|)^{-Np'}
\,\mathrm d\lambda_d(x)<\infty$.
Jeśli $p'=\infty$, mamy
$\|\varphi\|_\infty=q_{0,0}(\varphi)$.
W obu przypadkach otrzymujemy ciągłość.
W szczególności każda funkcja z $L^2$
może być traktowana jako dystrybucja temperowana,
co wykorzystamy przy porównywaniu transformat Fouriera.

\item Dla $a\in\R^d$ oraz $\alpha\in\N_0^d$ wzór
\[
 \langle\partial^\alpha\delta_a,\varphi\rangle
 =(-1)^{|\alpha|}\partial^\alpha\varphi(a)
\]
określa ciągłe przedłużenie pochodnej delty
na $\mathcal S$. Istotnie, moduł prawej strony
jest ograniczony przez $q_{0,|\alpha|}(\varphi)$.
Przypadek $\alpha=0$ daje samą deltę.

\item Każda dystrybucja $T\in\mathcal D'(\R^d)$
o zwartym nośniku $K$ jest temperowana.
Wybieramy $R\ge1$ takie, że $K\subset B_R(0)$
i funkcję odcinającą $\psi_R(x)=\psi(x/R)$
z dowodu twierdzenia~\ref{thm:pde8-density}.
Jest ona równa $1$ na otoczeniu $K$.
Dla $\varphi\in\mathcal S$ określamy
\[
 \langle\widetilde T,\varphi\rangle
 =\langle T,\psi_R\varphi\rangle.
\]
Iloczyn $\psi_R\varphi$ jest gładki
i ma nośnik w ustalonym zbiorze zwartym
$K_1=\overline{B_{2R}(0)}$.
Lokalne oszacowanie z rozdziału
~\ref{chap:pde-distributions} daje pewne
$C>0$ i $m\in\N_0$ takie, że
\[
 \begin{aligned}
 |\langle\widetilde T,\varphi\rangle|
 &\le C\max_{|\beta|\le m}\sup_{x\in K_1}
 |\partial^\beta(\psi_R\varphi)(x)|\\
 &\le C\,2^m q_{0,m}(\psi_R)q_{0,m}(\varphi).
 \end{aligned}
\]
Druga nierówność wynika z wielowskaźnikowej
reguły Leibniza. Wobec tego $\widetilde T$
jest ciągły na $\mathcal S$.
Jeśli $\varphi\in C_c^\infty$, to
$(1-\psi_R)\varphi$ jest testem o nośniku
rozłącznym z $K$. Z definicji nośnika
dystrybucji (sekcja~\ref{sec:pde7-support})
wynika $\langle T,(1-\psi_R)\varphi\rangle=0$.
Zatem $\widetilde T$ rzeczywiście przedłuża $T$.
Nie zależy od wyboru funkcji odcinającej,
ponieważ ciągłe przedłużenie jest jednoznaczne.

\item Dystrybucja regularna wyznaczona przez
$f(x)=e^{|x|^2}$ na $\R^d$ nie jest temperowana.
Funkcja $f$ jest lokalnie całkowalna, więc
określa dystrybucję na $C_c^\infty$.
Załóżmy, że ma ona ciągłe przedłużenie
$\widetilde T$ na $\mathcal S$.
Dla $R\ge1$ bierzemy
$\varphi_R(x)=\psi_R(x)e^{-|x|^2/2}$,
gdzie $\psi_R$ jest funkcją z dowodu
twierdzenia~\ref{thm:pde8-density}.
Funkcja $e^{-|x|^2/2}$ należy do $\mathcal S$,
a konstrukcja aproksymacji w tym dowodzie
daje $\varphi_R\to e^{-|x|^2/2}$ w $\mathcal S$.
Z ciągłości $\widetilde T$ wynikałaby
zbieżność $\widetilde T(\varphi_R)$
do skończonej liczby. Tymczasem
$0\le\psi_R\le1$, $\psi_R=1$ na $B_R(0)$
i $\varphi_R\in C_c^\infty$, więc
\[
 \begin{aligned}
 \langle\widetilde T,\varphi_R\rangle
 &=\int_{\R^d}\psi_R(x)e^{|x|^2/2}
 \,\mathrm d\lambda_d(x)\\
 &\ge\lambda_d(B_R(0))
 =R^d\lambda_d(B_1(0))\longrightarrow+\infty.
 \end{aligned}
\]
Jest to sprzeczność.
\end{enumerate}
\end{przyklad}

\section{Transformata Fouriera dystrybucji temperowanych}
\label{sec:pde8-fourier-distributions}

\begin{definicja}[Transformata na $\mathcal S'$]
\label{def:pde8-distFT}
Dla $T\in\mathcal S'(\R^d)$ definiujemy $\widehat T\in\mathcal S'(\R^d)$
przez \emph{transpozycję} operatora $\mathcal F$:
\[
 \langle\widehat T,\varphi\rangle
 =\langle T,\widehat\varphi\rangle,
 \qquad\varphi\in\mathcal S(\R^d).
\]
Analogicznie definiujemy $\mathcal F^{-1}T$ przez
$\langle\mathcal F^{-1}T,\varphi\rangle
=\langle T,\mathcal F^{-1}\varphi\rangle$.
\end{definicja}

Definicja jest poprawna: z twierdzeń
~\ref{thm:pde8-rules} i~\ref{thm:pde8-inversion}
wynika ciągłość obu operatorów na $\mathcal S$,
więc ich złożenia z $T$ są ciągłymi funkcjonałami
liniowymi. Nie stosujemy sprzężenia zespolonego:
jest to transpozycja względem liniowego działania
dystrybucji, nie operator sprzężony w przestrzeni Hilberta.

\begin{stwierdzenie}[Zgodność trzech definicji transformaty]
Jeśli $f\in L^1(\R^d)$, to
$\mathcal F T_f=T_{\widehat f}$, gdzie po prawej
stronie używamy transformaty całkowej.
Jeśli $f\in L^2(\R^d)$, to ta sama równość
zachodzi dla transformaty z twierdzenia Plancherela.
\end{stwierdzenie}
\begin{proof}
Dla $f\in L^1$ i $\varphi\in\mathcal S$
całka podwójna z modułu
$f(x)e^{-ix\cdot\xi}\varphi(\xi)$ wynosi
$\|f\|_1\|\varphi\|_1<\infty$. Stąd
\[
 \begin{aligned}
 \langle\widehat{T_f},\varphi\rangle
 &=\int_{\R^d}f(x)\widehat\varphi(x)
 \,\mathrm d\lambda_d(x)\\
 &=\int_{\R^d}\widehat f(\xi)\varphi(\xi)
 \,\mathrm d\lambda_d(\xi).
 \end{aligned}
\]
Funkcja $\widehat f$ jest ograniczona, więc
wyznacza dystrybucję temperowaną.
Dla $f\in L^2$ wybieramy $f_n\in\mathcal S$
zbiegające do $f$ w $L^2$.
Z twierdzenia~\ref{thm:pde8-plancherel}
wynika $\widehat f_n\to\widehat f$ w $L^2$.
Dla każdego testu $\varphi$ obie różnice działań
szacujemy przez nierówność Cauchy'ego--Schwarza:
\[
 \begin{aligned}
 |\langle T_{f_n}-T_f,\widehat\varphi\rangle|
 &\le\|f_n-f\|_2\|\widehat\varphi\|_2,\\
 |\langle T_{\widehat f_n}-T_{\widehat f},\varphi\rangle|
 &\le\|\widehat f_n-\widehat f\|_2\|\varphi\|_2.
 \end{aligned}
\]
Prawie wszędzie określone funkcje $L^2$
wyznaczają funkcjonały temperowane na mocy
sekcji~\ref{sec:pde8-tempered}. Przechodząc
do granicy w udowodnionej równości dla $f_n$,
otrzymujemy tezę.
\end{proof}

Pochodne dystrybucji temperowanych definiujemy
przez $\langle\partial_j T,\varphi\rangle
=-\langle T,\partial_j\varphi\rangle$.
Są temperowane dzięki ciągłości różniczkowania
na $\mathcal S$. Ogólniej, jeśli $a\in C^\infty(\R^d)$
i dla każdej pochodnej istnieją
$C_\gamma>0$, $M_\gamma\in\N_0$ z
$|\partial^\gamma a(x)|\le C_\gamma(1+|x|)^{M_\gamma}$,
to mnożenie przez $a$ jest ciągłe na $\mathcal S$.
Dla $m\in\N_0$ bierzemy
$M=\max_{|\gamma|\le m} M_\gamma$ i
$C=\max_{|\gamma|\le m} C_\gamma$.
Z reguły Leibniza dostajemy
\[
 q_{N,m}(a\varphi)\le 2^m C q_{N+M,m}(\varphi).
\]
Definiujemy więc $aT$ przez
$\langle aT,\varphi\rangle=\langle T,a\varphi\rangle$.
Te definicje zgadzają się z rachunkiem
na $\mathcal D'$ po ograniczeniu do testów
o zwartym nośniku. W szczególności można
mnożyć przez wielomiany oraz funkcje Schwartza;
dowolna funkcja gładka nie musi zachowywać
przestrzeni $\mathcal S$.

\begin{twierdzenie}[Rachunek na $\mathcal S'$]
\label{thm:pde8-dist-rules}
Transformata $\mathcal F:\mathcal S'\to\mathcal S'$ jest
liniowym automorfizmem, którego odwrotnością jest
$\mathcal F^{-1}$. Dla $T\in\mathcal S'(\R^d)$ i
$\alpha\in\N_0^d$ zachodzi
\[
 \mathcal F(\partial^\alpha T)=(i\xi)^\alpha\widehat T,
 \qquad \mathcal F(x^\alpha T)=i^{|\alpha|}
 \partial_\xi^\alpha\widehat T.
\]
\end{twierdzenie}
\begin{proof}
Dla każdego testu $\varphi\in\mathcal S$ mamy
\[
 \begin{aligned}
 \langle\mathcal F^{-1}\mathcal FT,\varphi\rangle
 &=\langle T,\mathcal F\mathcal F^{-1}\varphi\rangle
 =\langle T,\varphi\rangle,\\
 \langle\mathcal F\mathcal F^{-1}T,\varphi\rangle
 &=\langle T,\mathcal F^{-1}\mathcal F\varphi\rangle
 =\langle T,\varphi\rangle.
 \end{aligned}
\]
Liniowość wynika z definicji. Oba operatory
zachowują zbieżność słabą*: dla $T_n\to T$
testujemy $T_n$ odpowiednio przez ustaloną
funkcję $\mathcal F\varphi$ lub $\mathcal F^{-1}\varphi$.
Zatem są również ciągłe w topologii słabej*,
ponieważ każdy funkcjonał współrzędny po złożeniu
z nimi jest funkcjonałem współrzędnym tej topologii.
Dla pierwszej równości i testu $\varphi$ liczymy
\[
 \langle\widehat{\partial_jT},\varphi\rangle
 =-\langle T,\partial_j\widehat\varphi\rangle
 =\langle T,\widehat{i\xi_j\varphi}\rangle
 =\langle i\xi_j\widehat T,\varphi\rangle.
\]
Dla drugiej tożsamości korzystamy z równości
$\mathcal F(\partial_{\xi_j}\varphi)(x)
=ix_j\widehat\varphi(x)$:
\[
 \begin{aligned}
 \langle\mathcal F(x_jT),\varphi\rangle
 &=\langle T,x_j\widehat\varphi\rangle\\
 &=-i\langle T,\mathcal F(\partial_{\xi_j}\varphi)\rangle\\
 &=-i\langle\widehat T,\partial_{\xi_j}\varphi\rangle
 =i\langle\partial_{\xi_j}\widehat T,\varphi\rangle.
 \end{aligned}
\]
Iteracja daje oba wzory dla wielowskaźników,
ponieważ pochodne i mnożenia współrzędniowe
odpowiedniego rodzaju komutują.
\end{proof}

\begin{przyklad}[Cztery transformaty podstawowe]
\label{ex:pde8-four}
W $\mathcal S'(\R^d)$ zachodzą wzory
\[
 \widehat{\delta_0}=1,\quad
 \widehat 1=(2\pi)^d\delta_0,\quad
 \widehat{\delta_a}(\xi)=e^{-ia\cdot\xi},\quad
 \mathcal F(e^{ia\cdot x})=(2\pi)^d\delta_a.
\]
Dla dowolnego $a\in\R^d$ i testu $\varphi$ mamy
\[
 \langle\widehat{\delta_a},\varphi\rangle
 =\widehat\varphi(a)
 =\int_{\R^d}e^{-ia\cdot\xi}\varphi(\xi)
 \,\mathrm d\lambda_d(\xi).
\]
Przy $a=0$ otrzymujemy $\widehat{\delta_0}=1$.
Natomiast wzór odwrócenia dla $\varphi$ daje
\[
 \langle\mathcal F(e^{ia\cdot x}),\varphi\rangle
 =\int_{\R^d}e^{ia\cdot x}\widehat\varphi(x)
 \,\mathrm d\lambda_d(x)
 =(2\pi)^d\varphi(a).
\]
W szczególności przy $a=0$ mamy
$\widehat1=(2\pi)^d\delta_0$. Obie całki
są bezwzględnie zbieżne, ponieważ
$\varphi,\widehat\varphi\in L^1$.
Równości po prawej należy rozumieć \emph{dystrybucyjnie},
a nie jako zwykłe całki bezwzględnie zbieżne.
\end{przyklad}

\begin{twierdzenie}[Splot dystrybucji temperowanej z funkcją Schwartza]
\label{thm:pde8-convol-temp}
Niech $T\in\mathcal S'(\R^d)$ i $\varphi\in\mathcal S(\R^d)$.
Funkcja $(T*\varphi)(x):=\langle T(y),\varphi(x-y)\rangle$
jest gładka, ma wzrost co najwyżej wielomianowy
(wraz z każdą ustaloną pochodną) i wyznacza dystrybucję
temperowaną. Ponadto
$\mathcal F(T*\varphi)=\widehat T\,\widehat\varphi$ w $\mathcal S'$.
\end{twierdzenie}
\begin{proof}
Niech $F_x(y)=\varphi(x-y)$.
Z nierówności $(1+|y|)\le(1+|x|)(1+|x-y|)$
wynika dla każdego wielowskaźnika $\alpha$
\[
 q_{N,m}(\partial_x^\alpha F_x)
 \le(1+|x|)^N q_{N,m+|\alpha|}(\varphi).
\]
W szczególności wszystkie te funkcje należą do
$\mathcal S$. Dla $|h|\le1$ i $h\ne0$ podstawowe
twierdzenie rachunku całkowego, zastosowane
dwukrotnie wzdłuż kierunku $e_j$, daje
\[
 \frac{\partial_x^\alpha F_{x+he_j}
       -\partial_x^\alpha F_x}{h}
 -\partial_x^{\alpha+e_j}F_x
 =h\int_0^1(1-t)\partial_x^{\alpha+2e_j}F_{x+the_j}
 \,\mathrm d\lambda_1(t).
\]
Po różniczkowaniu względem $y$ i oszacowaniu
wagą wynika stąd
\[
 q_{N,m}(\text{różnica po lewej})
 \le \tfrac{|h|}{2}(2+|x|)^N
 q_{N,m+|\alpha|+2}(\varphi).
\]
Analogiczne oszacowanie pierwszej różnicy
przez całkę z pierwszej pochodnej daje
ciągłość $x\mapsto\partial_x^\alpha F_x$
w każdej półnormie (dla ogólnego przesunięcia
sumujemy różnice w kolejnych współrzędnych).
Z ciągłości $T$ możemy więc różniczkować
indukcyjnie:
$\partial_x^\alpha(T*\varphi)(x)
=\langle T,\partial_x^\alpha F_x\rangle$.
Jeśli $C,N,m$ są stałymi z globalnego oszacowania
~\eqref{eq:pde8-growth} dla $T$, to
\[
 |\partial_x^\alpha(T*\varphi)(x)|
 \le C(1+|x|)^N q_{N,m+|\alpha|}(\varphi).
\]
Jest to gładka funkcja o wzroście wielomianowym
wraz z pochodnymi, a dla $\alpha=0$ oszacowanie
zapewnia temperowaność jej dystrybucji regularnej.
Dla dowodu tożsamości Fouriera ustalmy
$\psi\in\mathcal S(\R^d)$ i oznaczmy
$A_x(y)=F_x(y)\widehat\psi(x)$.
W każdej półnormie mamy oszacowanie
\[
 q_{N,m}(A_x)\le(1+|x|)^N
 q_{N,m}(\varphi)|\widehat\psi(x)|.
\]
Prawa strona jest całkowalna względem $x$,
ponieważ $\widehat\psi\in\mathcal S$.
Niech $Q_R=[-R,R]^d$ i
$B_R(y)=\int_{Q_R}A_x(y)\,\mathrm d\lambda_d(x)$.
Możemy różniczkować całkę względem $y$,
gdyż pochodne $\varphi$ są ograniczone, a
$\widehat\psi$ jest całkowalna. Otrzymane
półnormy $B_R$ są skończone przez powyższe
oszacowanie, więc $B_R\in\mathcal S$.
Odwzorowanie $x\mapsto A_x$ jest ciągłe w każdej
półnormie, jak wykazano dla $F_x$, i dlatego
jednostajnie ciągłe w tej półnormie na $Q_R$.
Dla podziału $Q_R$ na prostokąty o średnicy
co najwyżej $\delta$ błąd sumy Riemanna
w półnormie $q_{N,m}$ nie przekracza
$\lambda_d(Q_R)$ razy
\[
 \sup_{\substack{x,z\in Q_R\\|x-z|\le\delta}}
 q_{N,m}(A_x-A_z),
\]
co dąży do zera przy $\delta\downarrow0$.
Ciągłość $T$ pozwala przejść do granicy sum,
a po stronie skalarnej całkujemy funkcję ciągłą
$x\mapsto\langle T,A_x\rangle$. Otrzymujemy
$\langle T,B_R\rangle=\int_{Q_R}\langle T,A_x\rangle
\,\mathrm d\lambda_d(x)$.

Niech $B(y)=\int_{\R^d}A_x(y)\,\mathrm d\lambda_d(x)$.
Całka i każda jej pochodna względem $y$ istnieją
dzięki ograniczoności pochodnych $\varphi$
i całkowalności $\widehat\psi$. Oszacowanie wag
daje $B\in\mathcal S$ oraz
\[
 q_{N,m}(B-B_R)\le q_{N,m}(\varphi)
 \int_{\R^d\setminus Q_R}(1+|x|)^N
 |\widehat\psi(x)|\,\mathrm d\lambda_d(x)\to0.
\]
Zatem $B_R\to B$ w $\mathcal S$.
Całka skalarna po całej przestrzeni jest
bezwzględnie zbieżna przez oszacowanie
wzrostu $T*\varphi$ i szybki zanik $\widehat\psi$.
Przechodząc do granicy w równości na $Q_R$,
dostajemy
\[
 \langle\mathcal F(T*\varphi),\psi\rangle
 =\int_{\R^d}(T*\varphi)(x)\widehat\psi(x)
 \,\mathrm d\lambda_d(x)
 =\langle T,B\rangle.
\]
Po podstawieniu $z=x-y$ i użyciu wzoru
$\widehat\psi(x)=\int_{\R^d}e^{-ix\cdot\xi}\psi(\xi)
\,\mathrm d\lambda_d(\xi)$ liczymy, stosując Fubiniego
(całka podwójna z modułu dla ustalonego $y$
wynosi $\|\varphi\|_1\|\psi\|_1$),
\[
 \int_{\R^d}\varphi(x-y)\widehat\psi(x)
 \,\mathrm d\lambda_d(x)
 =\int_{\R^d}e^{-iy\cdot\xi}\widehat\varphi(\xi)
 \psi(\xi)\,\mathrm d\lambda_d(\xi)
 =\mathcal F(\widehat\varphi\,\psi)(y).
\]
Stąd $\langle\mathcal F(T*\varphi),\psi\rangle
=\langle T,\mathcal F(\widehat\varphi\psi)\rangle
=\langle\widehat T,\widehat\varphi\psi\rangle$,
czyli $\mathcal F(T*\varphi)=\widehat T\,\widehat\varphi$.
W szczególności nie pojawia się dodatkowy czynnik
$(2\pi)^{-d}$: w naszej konwencji dotyczy on
transformaty \emph{iloczynu} funkcji, a nie splotu.
\end{proof}

\section{Mnożniki Fouriera: kiedy operator jest ograniczony?}
\label{sec:pde8-multipliers}

\begin{twierdzenie}[Mnożnik w $L^2$]
\label{thm:pde8-multiplier}
Niech $m:\R^d\to\C$ będzie funkcją mierzalną i
istotnie ograniczoną. Dla $f\in L^2(\R^d)$ istnieje dokładnie
jedna funkcja $T_mf\in L^2(\R^d)$ spełniająca
$\widehat{T_mf}=m\widehat f$ w sensie $L^2$.
Operator $T_m$ jest liniowy i ograniczony oraz
$\|T_m\|_{L^2\to L^2}=\|m\|_{L^\infty}$.
\end{twierdzenie}
\begin{proof}
Mnożenie przez $m$ jest ograniczone na $L^2$, a transformata
Fouriera jest unitarnym operatorem z dokładnością do
stałego czynnika, co daje istnienie, jednoznaczność i
oszacowanie z góry. Dla dolnego oszacowania bierzemy
$\widehat f$ proporcjonalną do funkcji charakterystycznej
zbioru dodatniej i skończonej miary, na którym
$|m|>\|m\|_\infty-\varepsilon$. Taki zbiór istnieje
(po przecięciu z dostatecznie dużą kulą), a po
$\varepsilon\downarrow0$ otrzymujemy równość norm.
\end{proof}

\begin{przyklad}[Transformata Riesza w $L^2$]
Dla $1\le j\le d$ połóżmy
$m_j(\xi)=-i\xi_j/|\xi|$ dla $\xi\ne0$ i $m_j(0)=0$.
Poprzednie twierdzenie definiuje ograniczony operator
$R_j=T_{m_j}:L^2\to L^2$, tzw. transformatę Riesza,
z $\|R_j\|\le1$. Wartość mnożnika w pojedynczym punkcie
nie zmienia operatora $L^2$. Ograniczoność w $L^p$ dla
$p\ne2$ jest istotnie trudniejszym rezultatem teorii
Calderóna--Zygmunda i nie wynika z samego warunku
$m_j\in L^\infty$.
\end{przyklad}

\section{Zastosowanie I: równanie ciepła i równanie falowe}
\label{sec:pde8-evolution}

\begin{twierdzenie}[Rozwiązania ewolucyjne poprzez transformatę]
\label{thm:pde8-evolution}
Niech $\kappa,c>0$ i $f,g\in\mathcal S(\R^d)$.
Wtedy funkcje
\begin{align*}
 u(t,x)&=\mathcal F^{-1}
 \bigl(e^{-\kappa t|\xi|^2}\widehat f(\xi)\bigr)(x),&&t\ge0,\\
 v(t,x)&=\mathcal F^{-1}\biggl(
 \cos(ct|\xi|)\widehat f(\xi)
 +\frac{\sin(ct|\xi|)}{c|\xi|}\widehat g(\xi)
 \biggr)(x),&&t\in\R,
\end{align*}
gdzie iloraz w $\xi=0$ definiujemy przez wartość graniczną $t$,
są gładkimi rozwiązaniami odpowiednio zagadnień
$u_t-\kappa\Delta u=0$, $u(0,\cdot)=f$ oraz
$v_{tt}-c^2\Delta v=0$, $v(0,\cdot)=f$, $v_t(0,\cdot)=g$.
W każdym ustalonym czasie funkcje przestrzenne należą do $\mathcal S$.
\end{twierdzenie}
\begin{proof}
Z reguły różniczkowania otrzymujemy po transformacji
$\partial_t\widehat u=-\kappa|\xi|^2\widehat u$ oraz
$\partial_{tt}\widehat v=-c^2|\xi|^2\widehat v$.
Są to dla każdego $\xi$ zwyczajne równania różniczkowe
z wskazanymi danymi początkowymi. Ich rozwiązania mają
właśnie zapis z tezy. Mnożniki są gładkie w $\xi$:
zwłaszcza $\sin(ct|\xi|)/(c|\xi|)$ ma rozwinięcie w szereg
potęg $|\xi|^2$ w otoczeniu zera.
Dla $t$ w dowolnym ograniczonym przedziale ich
pochodne po $t$ i $\xi$ mają co najwyżej wielomianowy wzrost,
stąd mnożą $\mathcal S$ w siebie. Uzasadnia to wszystkie
różniczkowania i stosowanie odwrotnej transformaty.
Wartości początkowe otrzymujemy przez podstawienie $t=0$.
\end{proof}

\begin{uwaga}[Dwie różne skale propagacji]
Równanie ciepła działa w częstotliwościach przez
$m_t(\xi)=e^{-\kappa t|\xi|^2}$: dla $t>0$ tłumi wysokie
częstotliwości bardzo szybko. Równanie falowe używa
$\cos(ct|\xi|)$ i $\sin(ct|\xi|)/(c|\xi|)$; nie ma
analogicznego wygładzania wszystkich wysokich częstotliwości.
To jest spektralna wersja zjawisk z rozdziałów
\ref{chap:pde-heat} i~\ref{chap:pde-wave}.
\end{uwaga}

\section[Potencjały Bessela i przestrzenie Sobolewa]{Potencjały Bessela i pierwsze przestrzenie $H^s$}
\label{sec:pde8-bessel}

\begin{definicja}[Przestrzeń Sobolewa typu Bessela]
\label{def:pde8-Hs}
Dla dowolnego $s\in\R$ definiujemy
\[
 H^s(\R^d)=\bigl\{u\in\mathcal S'(\R^d):
 (1+|\xi|^2)^{s/2}\widehat u\in L^2(\R^d)\bigr\}.
\]
Zapis przynależności do $L^2$ oznacza tutaj, że odpowiednia
dystrybucja jest reprezentowana przez funkcję kwadratowalną.
Normę tej przestrzeni definiujemy wzorem
$\|u\|_{H^s}^2=(2\pi)^{-d}\int_{\R^d}
(1+|\xi|^2)^s|\widehat u(\xi)|^2\,\mathrm d\lambda_d(\xi)$.
W tym wzorze $\widehat u$ jest reprezentowana przez funkcję
lokalnie $L^2$: otrzymujemy ją mnożąc funkcję $L^2$
przez $(1+|\xi|^2)^{-s/2}$.
Operator $J^s=(I-\Delta)^{s/2}$ określamy na $\mathcal S'$
przez mnożnik $(1+|\xi|^2)^{s/2}$.
Dla ujemnego $s$ element $H^s$ może być autentyczną
dystrybucją, a nie funkcją lokalnie całkowalną;
wymóg $L^2$ odnosi się do ważonej transformaty,
nie do samej dystrybucji $u$.
\end{definicja}

Dla $m\in\N_0$ zapis $W^{m,2}(\R^d)$ oznacza przestrzeń
funkcji $u\in L^2(\R^d)$, których wszystkie pochodne słabe
$\partial^\alpha u$, $|\alpha|\le m$, należą do $L^2$.

\begin{uwaga}[Szeregi Fouriera na torusie a transformata Fouriera na $\R^d$]
Rozdział~\ref{chap:pde-fourier} pracował na torusie lub
na odcinku po okresowym przedłużeniu i korzystał z rozwinięć
w szeregi Fouriera, czyli dyskretnego spektrum częstości.
Tutaj pracujemy na całej przestrzeni $\R^d$, więc częstość
$\xi$ zmienia się ciągle i pojawia się całka po $\R^d$.
Obie teorie są zgodne pojęciowo, ale różnią się nośnikiem,
normalizacją oraz tym, że w części okresowej integrujemy na
$\T^d$ względem ustalonej miary, a tutaj na $\R^d$ względem
miary Lebesgue'a.
\end{uwaga}

\begin{twierdzenie}[Podstawowe własności $H^s$]
\label{thm:pde8-Hs}
Dla wszystkich $s,t\in\R$, $s\ge t$, zachodzą:
\begin{enumerate}
\item $H^s(\R^d)$ jest przestrzenią Hilberta, a
 $J^s:H^s\to L^2$ jest izometrią na całą przestrzeń $L^2$.
\item $H^s\hookrightarrow H^t$ ciągle i
 $\|u\|_{H^t}\le\|u\|_{H^s}$.
\item Dla $\alpha\in\N_0^d$, $\partial^\alpha:H^s\to
 H^{s-|\alpha|}$ jest ograniczony.
\item Dla $m\in\N_0$ zachodzi $H^m(\R^d)=W^{m,2}(\R^d)$
 ze wzajemnie równoważnymi normami.
\end{enumerate}
\end{twierdzenie}
\begin{proof}
Pierwsza teza wynika z utożsamienia $H^s$ poprzez
$J^s$ z $L^2$ i odwracalności mnożnika
$(1+|\xi|^2)^{s/2}$ na $\mathcal S'$. Pozostałe
oszacowania opierają się na nierównościach
$(1+|\xi|^2)^t\le(1+|\xi|^2)^s$ i
$|\xi^\alpha|^2\le(1+|\xi|^2)^{|\alpha|}$.
Gdy $m$ jest całkowite nieujemne, wielomian
$(1+|\xi|^2)^m$ jest równoważny
$\sum_{|\alpha|\le m}|\xi^\alpha|^2$,
ze stałymi zależnymi tylko od $m,d$. Dzięki
Plancherelowi warunek $u\in H^m$ jest równoważny
istnieniu wszystkich pochodnych dystrybucyjnych
$\partial^\alpha u\in L^2$ dla $|\alpha|\le m$.
W przypadku $m=0$ warunek ten mówi po prostu $u\in L^2$.
\end{proof}

\begin{wniosek}[Kiedy delta należy do ujemnej przestrzeni Sobolewa?]
\label{cor:pde8-delta-Hs}
Dla $s\in\R$ zachodzi
\[
 \delta_0\in H^s(\R^d)\quad\Longleftrightarrow\quad
 s<-\frac d2.
\]
\end{wniosek}
\begin{proof}
Ponieważ $\widehat{\delta_0}=1$, warunkiem jest
$\int_{\R^d}(1+|\xi|^2)^s\,\mathrm d\lambda_d(\xi)<\infty$.
Całka po kuli jednostkowej jest skończona. Poza nią,
przechodząc do współrzędnych biegunowych, porównujemy
całkę z $\int_1^\infty r^{2s+d-1}\,\mathrm d\lambda_1(r)$,
która jest skończona dokładnie wtedy, gdy $2s+d<0$.
\end{proof}

\section{Zastosowanie II: eliptyczne równanie z parametrem}
\label{sec:pde8-elliptic}

\begin{twierdzenie}[Równanie Yukawy--Helmholtza]
\label{thm:pde8-helmholtz}
Niech $a>0$ i $f\in L^2(\R^d)$. Istnieje dokładnie jedno
$u\in H^2(\R^d)$ takie, że $(-\Delta+a^2)u=f$
w sensie dystrybucji. Zachodzi
\[
 \widehat u(\xi)=\frac{\widehat f(\xi)}{|\xi|^2+a^2},
 \qquad \|u\|_{H^2}\le C_a\|f\|_{L^2},
\]
gdzie $H^2$ zdefiniowano w~\ref{def:pde8-Hs},
a $C_a=\max\{1,a^{-2}\}$ dla przyjętej tam normy.
\end{twierdzenie}
\begin{proof}
Funkcja $(|\xi|^2+a^2)^{-1}$ jest ograniczona.
Ponadto $(1+|\xi|^2)/(|\xi|^2+a^2)\le\max\{1,a^{-2}\}$.
Zatem zadany wzór wyznacza $u\in H^2$ i podane
oszacowanie przez twierdzenie Plancherela. Po
pomnożeniu przez $|\xi|^2+a^2$ otrzymujemy równanie.
Jeżeli $(-\Delta+a^2)u=0$ dla $u\in H^2$, to
$(|\xi|^2+a^2)\widehat u=0$ prawie wszędzie,
a więc $u=0$. To daje jednoznaczność.
\end{proof}

\begin{uwaga}[Dlaczego przypadek $a=0$ jest inny?]
Dla czystego równania Poissona mnożnik formalny
$|\xi|^{-2}$ ma osobliwość w zerze. Nie można więc
bez dodatkowej analizy definiować dla dowolnego
$f\in L^2$ rozwiązania $u\in H^2$ przez ten wzór.
Dokładniej, dla $1\le d\le4$ i $f\in\mathcal S(\R^d)$
z $\widehat f(0)\ne0$ iloraz
$|\xi|^{-2}\widehat f(\xi)$ nie jest funkcją $L^2$
w pobliżu zera, gdyż jego kwadrat daje po przejściu
do współrzędnych biegunowych całkę porównywalną z
$\int_0^1r^{d-5}\,\mathrm d\lambda_1(r)=\infty$.
W wymiarach $d\ge5$ ta konkretna przeszkoda
niskoczęstotliwościowa znika dla takiego gładkiego $f$;
nie zmienia to faktu, że dzielenie przez $|\xi|^2$
nie jest ograniczonym mnożnikiem $L^2\to H^2$
bez dodatkowych warunków na dane. Dodatkowo funkcje harmoniczne tworzą jądro
operatora $-\Delta$ w szerszych klasach rozwiązań.
Nie oznacza to braku rozwiązań dystrybucyjnych;
rozdział~\ref{chap:pde-distributions} konstruuje je
z użyciem funkcji fundamentalnej.
\end{uwaga}

\section{Zasada nieoznaczoności jako konsekwencja Plancherela}
\label{sec:pde8-uncertainty}

\begin{twierdzenie}[Nierówność Heisenberga w jednej współrzędnej]
\label{thm:pde8-uncertainty}
Jeśli $0\ne f\in\mathcal S(\R^d)$ i $1\le j\le d$, to
\[
 \|x_jf\|_{L^2}\,\|\xi_j\widehat f\|_{L^2}
 \ge\frac12\,\|f\|_{L^2}\,\|\widehat f\|_{L^2}.
\]
\end{twierdzenie}
\begin{proof}
Całkując pochodną $\partial_j(x_j|f|^2)$ i używając
zaniku funkcji Schwartza, dostajemy
$\|f\|_2^2=-2\operatorname{Re}\int_{\R^d}x_j\overline f\,
\partial_jf\,\mathrm d\lambda_d(x)$.
Nierówność Cauchy'ego--Schwarza daje
$\|x_jf\|_2\|\partial_jf\|_2\ge\tfrac12\|f\|_2^2$.
Z twierdzenia Plancherela i
$\widehat{\partial_jf}=i\xi_j\widehat f$ wynika
$\|\partial_jf\|_2=(2\pi)^{-d/2}\|\xi_j\widehat f\|_2$
oraz $\|f\|_2=(2\pi)^{-d/2}\|\widehat f\|_2$.
Po podstawieniu otrzymujemy tezę.
\end{proof}

\section{Podsumowanie i zakres dalszych zagadnień}
Transformata Fouriera jest automorfizmem przestrzeni Schwartza i
jej przestrzeni dualnej. Jest podstawowym narzędziem badania wielu
liniowych PDE o stałych współczynnikach, lecz sama nie zapewnia
rozwiązywalności ani jednoznaczności w dowolnej klasie funkcji;
nie rozwiązuje też automatycznie problemów brzegowych ani równań
ze zmiennymi współczynnikami. Mnożniki ograniczone działają na $L^2$,
lecz dla $L^p$ z $p\ne2$ potrzebna jest głębsza teoria.
W następnych rozdziałach zbudujemy systematycznie
przestrzenie Sobolewa na obszarach i ich twierdzenia o zanurzeniach.
W ten sposób domkniemy łańcuch pojęciowy rozpoczęty w
rozdziale~\ref{chap:pde-distributions}: od słabych pochodnych,
przez narzędzia Fouriera na całej przestrzeni, do funkcjonałowej
teorii $W^{k,p}(\Omega)$ potrzebnej dla metod wariacyjnych i
zagadnień brzegowych.



